% ============== CHAPITRE 2 — ANALYSE ET SPÉCIFICATION ================
\chapter{Analyse et spécification des besoins}

\section*{Introduction}
Ce chapitre est consacré à l'\textbf{analyse} du système à réaliser. Il identifie
d'abord les acteurs qui interagissent avec l'application, spécifie ensuite les
besoins fonctionnels et non fonctionnels, présente le diagramme de cas
d'utilisation global, puis détaille par une description textuelle les principaux
cas d'utilisation, regroupés par module fonctionnel.

% ----------------------------------------------------------------------
\section{Identification des acteurs}
Trois profils interagissent avec l'application, conformément aux rôles gérés par
le module de sécurité (table \texttt{SECURITY\_USERS}). Le
tableau~\ref{tab:acteurs} en précise les responsabilités.

\begin{table}[htbp]
\centering
\caption{Identification des acteurs}
\label{tab:acteurs}
\renewcommand{\arraystretch}{1.25}
\begin{tabularx}{\textwidth}{>{\bfseries}p{4cm}X}
\toprule
\textmd{\textbf{Acteur}} & \textbf{Responsabilités}\\
\midrule
Administrateur (\texttt{ADMIN}) & Gère les agences, les employés et les clients ;
valide les demandes de création de clients et d'employés ; consulte le journal
d'audit ; supervise le volet décisionnel.\\
Directeur commercial (\texttt{DIRECTEUR\_COMMERCIAL}) & Soumet les demandes de
création (clients, employés) ; valide les demandes de modification des données
des employés ; consulte les objectifs et les indicateurs.\\
Utilisateur (\texttt{USER}) & Consulte les données autorisées, effectue son
pointage et soumet des demandes de modification de ses propres informations.\\
\bottomrule
\end{tabularx}
\end{table}

% ----------------------------------------------------------------------
\section{Spécification des besoins}

\subsection{Besoins fonctionnels}
Le système doit permettre :
\begin{itemize}[noitemsep]
  \item l'authentification sécurisée et la gestion des rôles et des droits
        d'accès (\texttt{ADMIN}, \texttt{DIRECTEUR\_COMMERCIAL}, \texttt{USER}) ;
  \item la gestion (CRUD) des agences, des employés et des clients ;
  \item la consultation des objectifs commerciaux et le suivi de la production ;
  \item le \textbf{pointage (présence) des employés} : saisie de l'heure d'arrivée
        et de départ, statut (présent / retard / absent), en mode automatique
        (bouton) ou manuel (administrateur) ;
  \item la soumission et la validation des demandes (clients, employés,
        modifications) selon le cycle \texttt{EN\_ATTENTE} $\rightarrow$
        \texttt{EN\_COURS} $\rightarrow$ \texttt{EFFECTUE}/\texttt{REFUSE} ;
  \item la notification des demandes en attente et la journalisation des actions ;
  \item la consolidation des données par une chaîne \textbf{ETL} et la restitution
        d'indicateurs via un tableau de bord et Power BI ;
  \item la \textbf{segmentation des agences} par apprentissage non supervisé.
\end{itemize}

\subsection{Besoins non fonctionnels}
Le tableau~\ref{tab:bnf} récapitule les besoins non fonctionnels.

\begin{table}[htbp]
\centering
\caption{Besoins non fonctionnels}
\label{tab:bnf}
\renewcommand{\arraystretch}{1.2}
\begin{tabularx}{\textwidth}{>{\bfseries}p{3.4cm}X}
\toprule
\textmd{\textbf{Besoin}} & \textbf{Description}\\
\midrule
Sécurité & Authentification, gestion des rôles et protection des opérations
sensibles par des workflows de validation.\\
Fiabilité & Intégrité des données garantie par le SGBD et les transactions JTA.\\
Performance & Temps de réponse acceptable sur les volumes manipulés.\\
Ergonomie & Interface claire et cohérente, navigation par menu latéral.\\
Maintenabilité & Architecture en couches, code modulaire et faiblement couplé.\\
Portabilité & Application Jakarta EE standard, déployable sur serveur compatible.\\
\bottomrule
\end{tabularx}
\end{table}

% ----------------------------------------------------------------------
\section{Diagramme de cas d'utilisation global}
La figure~\ref{fig:uc_global} présente le \textbf{diagramme de cas d'utilisation
global}, qui représente d'un point de vue fonctionnel les interactions entre les
acteurs et les services offerts par le système.

\begin{figure}[p]
  \centering
  \rotatebox{90}{%
    \includegraphics[width=0.92\textheight,height=\textwidth,keepaspectratio]{uc_global.png}}
  \caption{Diagramme de cas d'utilisation global}
  \label{fig:uc_global}
\end{figure}

On y distingue les trois acteurs et leurs cas d'utilisation respectifs. Le cas
\emph{S'authentifier} est commun à tous ; il est relié aux autres par une relation
\texttt{\guillemotleft include\guillemotright}, l'accès aux fonctionnalités exigeant une
authentification préalable. L'\textbf{administrateur} concentre les cas de gestion
et de validation ; le \textbf{directeur commercial} soumet les demandes et consulte
les indicateurs ; l'\textbf{utilisateur} se limite à la consultation, au pointage
et aux demandes le concernant.

% ----------------------------------------------------------------------
\section{Description textuelle des cas d'utilisation}
Cette section détaille, sous forme de fiches, les principaux cas d'utilisation,
regroupés par module fonctionnel.

\subsection{Authentification et sécurité}

\cudesc{tab:cu_auth}{S'authentifier}
{Permettre à l'utilisateur de se connecter à l'application avec son identifiant et
son mot de passe.}
{Administrateur / Directeur commercial / Utilisateur}
{L'utilisateur doit disposer d'un compte actif dans la table \texttt{SECURITY\_USERS}.}
{Une session est ouverte et l'utilisateur est redirigé vers le tableau de bord
correspondant à son rôle.}
{(1) L'utilisateur saisit son identifiant et son mot de passe ; (2) le
\texttt{SecurityService} vérifie les informations sur la table de sécurité ; (3) en
cas de succès, une session est ouverte et l'utilisateur est redirigé ; (4) sinon,
un message d'erreur est affiché.}

\cudesc{tab:cu_roles}{Gérer les rôles et les droits d'accès}
{Permettre à l'administrateur d'attribuer un rôle à chaque utilisateur afin de
contrôler l'accès aux fonctionnalités.}
{Administrateur}
{L'administrateur doit être authentifié avec le rôle \texttt{ADMIN}.}
{Le rôle est enregistré et les droits d'accès de l'utilisateur sont mis à jour.}
{(1) L'administrateur consulte la liste des utilisateurs ; (2) il sélectionne un
utilisateur ; (3) il lui attribue un rôle (\texttt{ADMIN},
\texttt{DIRECTEUR\_COMMERCIAL} ou \texttt{USER}) ; (4) le système applique le
contrôle d'accès correspondant.}

\subsection{Gestion du référentiel : agences et employés}

\cudesc{tab:cu_agence}{Gérer les agences}
{Permettre à l'administrateur de créer, consulter, modifier et supprimer les
agences du réseau.}
{Administrateur}
{L'administrateur doit être authentifié avec le rôle \texttt{ADMIN}.}
{L'agence créée ou modifiée est enregistrée dans la base et la liste est mise à
jour.}
{(1) L'administrateur accède à l'interface de gestion des agences ; (2) il
sélectionne une action (créer, modifier, supprimer) ; (3) il saisit ou met à jour
les informations de l'agence ; (4) l'\texttt{AgenceService} persiste la
modification et la liste est rafraîchie.}

\cudesc{tab:cu_employe}{Gérer les employés}
{Permettre à l'administrateur de gérer les employés et de les rattacher à leur
agence.}
{Administrateur}
{L'administrateur doit être authentifié et l'agence de rattachement doit exister.}
{L'employé est enregistré et rattaché à son agence.}
{(1) L'administrateur accède à l'interface de gestion des employés ; (2) il saisit
les informations de l'employé ; (3) il sélectionne l'agence de rattachement ; (4)
l'\texttt{UtilisateurService} enregistre l'employé associé à l'agence.}

\subsection{Gestion commerciale : clients et objectifs}

\cudesc{tab:cu_client}{Gérer les clients}
{Permettre à l'administrateur de gérer le portefeuille clients et d'affecter
chaque client à un gestionnaire.}
{Administrateur}
{L'administrateur doit être authentifié et le gestionnaire de rattachement doit
exister.}
{Le client est enregistré et associé à son gestionnaire dans la base.}
{(1) L'administrateur accède à l'interface de gestion des clients ; (2) il saisit
ou met à jour les informations du client ; (3) il sélectionne le gestionnaire
responsable ; (4) le \texttt{ClientService} enregistre le client et met à jour le
portefeuille.}

\cudesc{tab:cu_objectif}{Suivre les objectifs commerciaux}
{Permettre au directeur commercial de consulter les objectifs et la production par
agence.}
{Directeur commercial / Administrateur}
{L'acteur doit être authentifié et des objectifs doivent être renseignés pour la
période.}
{Les indicateurs de production (comptes, crédits, épargne) sont restitués par
agence.}
{(1) L'acteur accède à l'interface des objectifs ; (2) il sélectionne une agence
et une période ; (3) l'\texttt{ObjectifService} agrège la production ; (4) les
indicateurs sont affichés.}

\subsection{Pointage et workflows de demandes}

\cudesc{tab:cu_pointage}{Pointer la présence}
{Permettre à l'employé d'enregistrer son arrivée et son départ, avec
détermination automatique du statut.}
{Utilisateur (employé)}
{L'employé doit être authentifié.}
{Le pointage du jour est enregistré avec son statut (\texttt{PRESENT} ou
\texttt{RETARD}).}
{(1) L'employé clique sur « Pointer » ; (2) le \texttt{PointageService} relève
l'heure courante et recherche le pointage du jour ; (3) s'il n'existe pas, il
détermine le statut (\texttt{RETARD} au-delà de 9 h) et l'enregistre ; (4) sinon,
il met à jour l'enregistrement existant.}

\cudesc{tab:cu_demande}{Valider une demande}
{Permettre à l'administrateur d'approuver ou de refuser une demande soumise.}
{Administrateur}
{Une demande au statut \texttt{EN\_ATTENTE} doit exister.}
{Le statut de la demande évolue vers \texttt{EFFECTUE} ou \texttt{REFUSE} et
l'action est journalisée.}
{(1) L'administrateur consulte les demandes en attente ; (2) il sélectionne une
demande ; (3) il l'approuve ou la refuse avec un commentaire ; (4) le service met
à jour le statut, applique la modification le cas échéant et journalise l'action.}

\subsection{Volet décisionnel}

\cudesc{tab:cu_etl}{Alimenter le datamart (ETL)}
{Consolider les données opérationnelles en un datamart agrégé par agence.}
{Processus ETL / Administrateur}
{Les données sources (agences, employés, clients, objectifs, pointage) doivent
être disponibles.}
{Le datamart en étoile est produit et mis à la disposition de la restitution et de
la segmentation.}
{(1) L'étape \emph{Extract} lit les données sources ; (2) l'étape \emph{Transform}
les nettoie et les agrège par agence ; (3) l'étape \emph{Load} écrit le datamart
(\texttt{dim\_agence}, \texttt{fait\_agence}).}

\cudesc{tab:cu_dashboard}{Consulter le tableau de bord}
{Restituer aux responsables une vision synthétique de l'activité.}
{Administrateur / Directeur commercial}
{L'acteur doit être authentifié.}
{Les indicateurs clés (effectifs, présence, production) sont affichés.}
{(1) L'acteur se connecte ; (2) le \texttt{HomeServlet} calcule les statistiques ;
(3) la page \texttt{home.jsp} affiche les cartes d'indicateurs et les graphiques
ApexCharts.}

\cudesc{tab:cu_clustering}{Segmenter les agences}
{Regrouper automatiquement les agences en profils homogènes à partir de leurs
indicateurs d'activité.}
{Analyste décisionnel / Administrateur}
{Le datamart agrégé par agence (\texttt{fait\_agence}) doit être disponible.}
{Les agences sont réparties en segments cohérents, caractérisés et exploitables
pour la décision.}
{(1) L'analyste lance la segmentation ; (2) le script standardise les variables ;
(3) il détermine le nombre de clusters puis compare les modèles ; (4) il retient
le meilleur modèle, projette en PCA et restitue le profil de chaque segment.}

\section*{Conclusion}
Ce chapitre a identifié les acteurs, spécifié les besoins fonctionnels et non
fonctionnels et formalisé les cas d'utilisation du système. Le chapitre suivant
traduit cette analyse en une \textbf{conception} détaillée : architecture,
modèle de classes et scénarios dynamiques.
